\section{Introduction}
\label{sec:intro}

Auxiliary losses that pull a network's features toward a known concept
direction are a standard way to inject prior knowledge. We report a
regime in which they systematically backfire. Take an EEG emotion
classifier that has already converged, and add a loss aligning its
features with a valence direction: across twenty-five such recipes on
FACED~\citep{chen2023faced}, none improve balanced accuracy and fifteen
significantly reduce it. The same loss that helps a weak backbone hurts
a strong one. We call this the \emph{saturation regularity}: once task
labels alone have driven the network onto the target concept, further
supervision can only deform an already-converged representation, without
improving generalization. The finding is doubly useful. It explains a
widespread failure mode of concept-alignment supervision, and, by
locating where the recoverable gain does live---in the seed-specific
residual that supervision cannot exploit---it yields a new state of the
art through ensembling rather than supervising.

To make ``the concept the network has saturated on'' concrete, we
extract it explicitly. Pass nine short emotion-evocative sentences
through a language model, average the late-layer hidden states class by
class, and take the first principal component; the resulting
\emph{valence axis} (V-axis) predicts movie-review sentiment at AUC
$0.832$ zero-shot, tracks the cohort-averaged EEG response of $123$
subjects at the emotion-category level ($r{=}0.87$), and is rediscovered
without supervision by every reasonably-strong EEG classifier we tested.
Recent work shows that language models' internal features align with
human cognition across high-level
concepts~\citep{kim2018tcav,arditi2024refusal}; we use the V-axis as a
\emph{probe} of what brain-decoding networks encode, not as a training
signal---which, as the saturation regularity shows, a converged network
cannot benefit from.

\textbf{A nine-sentence axis.} The nine stories are one per
FACED~\citep{chen2023faced} class; the language model is Qwen3-4B,
chosen as the lead because it sits at the centre of the alignment
manifold in our $14$-LLM sweep, with the highest per-stimulus
alignment to a behavioural valence reference (the recipe itself is
model-agnostic, see Appendix~\ref{app:per-llm-deep-dive}). The SST-2
benchmark grounds the AUC $0.832$ headline:
the V-axis projection alone is within $0.005$ of a $5$k-example
supervised logistic regression~\citep{socher2013recursive}. And the
direction is not idiosyncratic to any one model: across the
$14$ language models from $560$M to $32$B parameters, the V-axis
converges on essentially the same direction (within Qwen3,
off-diagonal $r{=}{+}0.995$ over $15$ size-pairs).

\textbf{The same axis lives in the brain.} The $r{=}0.87$ above
unpacks as follows. A ridge regression on $160$ channel-band
features of the FACED cohort EEG predicts the V-axis position of
each of the $28$ stimuli; the regressor uses leave-one-stimulus-out
cross-validation, no per-subject fitting, no auxiliary supervision.
The Qwen3-4B V-axis itself reaches $r{=}0.80$ ($p<10^{-5}$) on this
ridge; a CLIP-text variant of the V-axis built from the same $28$
stimulus descriptions tightens this to $r{=}0.87$ ($p<10^{-9}$).
The third claim --- that EEG networks rediscover the V-axis without
ever being shown it --- is quantified across $36$ FACED-9 checkpoints
trained on the $9$ emotion labels alone: per-checkpoint V-axis
encoding strength in the class-mean subspace predicts balanced
accuracy at $r{=}{+}0.885$ ($p{=}7.8{\times}10^{-13}$,
Figure~\ref{fig:universal-vaxis}). The same valence direction is
recovered from inside language models and from inside human
electrophysiology, by independent procedures.

\begin{figure}[ht]
\centering
\IfFileExists{figures/landmark/lf1_universal_vaxis_hero.pdf}
  {\includegraphics[width=\linewidth]{landmark/lf1_universal_vaxis_hero.pdf}}
  {\includegraphics[width=\linewidth]{F1_universal_vaxis.pdf}}
\caption[V-axis hero]{\textbf{A single one-dimensional valence direction, extracted from
nine LLM emotion stories, transfers to text sentiment and to cohort
brain response.}
\textbf{(a)} Zero-shot probe AUC / $|r|$ of the V-axis (PC1 of nine
Qwen3-4B story-class centroids at the penultimate layer) on standard
text-sentiment benchmarks; no labels are seen during axis construction.
\textbf{(b)} Stimulus-level scatter of LLM V-axis projection vs.\
cohort EEG response on FACED ($n{=}28$ stimuli, $123$ subjects), a
fully nested permutation test placing the alignment beyond every
relabelling ($p{\le}2{\times}10^{-4}$); inset shows the
random-direction null ($r{=}0.07$). The same one-dimensional direction
explains text sentiment and predicts evoked EEG without sharing data
across the two settings.
}
\label{fig:universal-vaxis}
\end{figure}

\textbf{Yet supervision fails.} The natural way to use such an
alignment is an auxiliary loss pulling the EEG network's penultimate
features toward the V-axis. Across $25$ such recipes---knowledge
distillation~\citep{hinton2015distilling}, representational
similarity~\citep{kriegeskorte2008rsa},
contrastive~\citep{khosla2020supcon,vandenoord2018cpc}, PEFT
adapters~\citep{hu2022lora}, and topographic
losses~\citep{padakanti2025meg}---none help: fifteen produce
statistically significant decrements and zero produce gains, five
families showing \emph{monotonic destruction}. The effect flips sign
with baseline strength (no measurable effect at BACC $\le 0.62$, a
significant negative at BACC $\ge 0.66$). We call this
\emph{saturation}: distinct from overfitting, it is what happens once
task labels alone have driven the network onto the target direction,
so added supervision moves the class-mean component further along the
axis ($\Delta|r|$ up to ${+}0.36$) while test accuracy falls and
training accuracy stays at ceiling.

\textbf{Where the gain lives.} The same diagnosis locates the
recoverable gain in the within-class residual that supervision cannot
reach: the class-mean basin is shared across seeds, while the residual
is seed-specific and load-bearing. Ensembling across seeds averages
out residual noise without disturbing the basin---per-checkpoint
residual encoding predicts leave-one-out contribution at $r{=}{+}0.74$,
and a $10$-checkpoint ensemble reaches FACED-9 SOTA at
$\mathbf{0.6948}$ (val-selected single $\mathbf{0.6755}$), a $+10.5\%$
improvement over EMOD~\citep{chen2025emod}. Replacing the
language-derived KD prototypes with random orthonormal vectors changes
BACC by $\le 0.003$: the gain is the \emph{shape} of the loss, not the
\emph{content} of the prototypes, so the V-axis is a probe of what the
network saturated on, not the supervised signal.

\paragraph{Contributions.} \textbf{(1)} A nine-sentence \emph{V-axis
probe}: a one-dimensional valence direction extracted from $9$
language-model emotion stories that converges across $14$ LLMs,
predicts text sentiment zero-shot, predicts cohort EEG response, and
is rediscovered without supervision by $36$ EEG classifiers
(\S\ref{sec:vaxis-llms}--\S\ref{sec:topography}).
\textbf{(2)} The \emph{saturation regularity}: across $25$
representation-alignment recipes on FACED-9, $15$ produce significant
decrements and $0$ produce gains, with monotonic destruction in $5$
families and a sharp transition in the $[0.62, 0.66]$ BACC band; the
loss moves the class-mean component by $\Delta|r|{=}{+}0.01$ to
${+}0.36$ while training accuracy stays at ceiling and test accuracy
falls (\S\ref{sec:saturation}).
\textbf{(3)} A residual-diversity ensemble turning the mechanism into
a positive prescription: per-checkpoint residual encoding predicts
leave-one-out contribution at $r{=}{+}0.74$, and a $10$-checkpoint
ensemble reaches FACED-9 SOTA at $\mathbf{0.6948}$
(\S\ref{sec:convergence},~\S\ref{sec:sota}). The saturation
regularity and the cross-architecture V-axis convergence both
replicate on a second EEG dataset,
SEED-V~\citep{liu2022seedv} (Appendix~\ref{app:seedv-full}).

\paragraph{What this paper does not claim.} To keep the scope
precise, we state four non-claims up front. First, we do not claim
the V-axis is a privileged or uniquely useful direction: a generic
valence axis performs identically as a probe (\S\ref{sec:vaxis-llms}),
and inside the ensemble recipe the language-derived prototypes are
interchangeable with random orthonormal vectors ($\Delta\le0.003$
BACC). Second, we do not claim single-trial or within-subject brain
decoding: the brain correlation is measured on the EEG response
averaged across $123$ subjects at the emotion-category level, and the
within-subject effect is near zero ($r{=}0.03$). Third, we do not
claim that auxiliary alignment is universally harmful --- the same
valence loss \emph{helps} a weak backbone ($+0.057$ BACC on LaBraM);
saturation is specifically a property of the \emph{converged} regime.
Fourth, we do not present saturation as a proven theorem for all
architectures and datasets: we establish it empirically on FACED-9
and replicate the saturation regularity on SEED-V. Limitations are
collected in \S\ref{sec:discussion}.
